% GENERATED FILE. Edit canonical lessons, then run npm run generate:books.
% canonical-source-hash: fnv1a64:fb9c8a86e47c03dc
% canonical-lessons: ZH-C155-cong, ZH-W155-cong-from, ZH-C155-cong-dao, ZH-C155-chufa, ZH-C155-kaimen, ZH-C155-guanmen, ZH-C155-reading, ZH-C155-practice

\chapter{Trains and Opening Hours}
\label{ch:zh-trains-and-opening-hours}
\hlchaptermodality{155}

\begin{chapteropening}
\textbf{By the end of this chapter:} \emph{I can find times in a timetable and opening hours.}

Reading a station timetable and a bookshop's opening hours, with \zh{从} … \zh{到} …, setting off, opening and closing, built on one new written character: \zh{从}.
\end{chapteropening}

\section[\texorpdfstring{\zh{从} (cóng)}{从 (cóng)}]{\zh{从} (cóng) — from}

\label{lesson:ZH-C155-cong}



\begin{quote}
A new chapter, and it opens on the word for where things start.

\begin{itemize}
\raggedright
  \item \emph{Recall:} say \zh{你要水还是要饭} — \emph{nǐ yào shuǐ háishi yào fàn} — \textbf{R1}, the chapter you have finished
  \item \emph{Recall:} say \emph{bié}, don't — \textbf{R1}
  \item \emph{Recall:} say \emph{cónglái}, ever — \textbf{R2}, a few chapters back
\end{itemize}
\end{quote}

\subsection*{You'll want to know: \zh{从}}
\textbf{\zh{从}} — \emph{cóng} — \textquotedblleft{}from\textquotedblright{}.

\begin{quote}
\textbf{\zh{从这儿到大学}} \emph{cóng zhèr dào dàxué} — from here to the university
\end{quote}

\begin{quote}
\textbf{\zh{我从中国来}} \emph{wǒ cóng Zhōngguó lái} — I come from China.
\end{quote}

The next lesson puts the character in your hand. It is the cóng of \emph{cónglái}.

\begin{cousinweb}
\textbf{\zh{从}} (\emph{cóng}) is inherited from Middle Chinese \emph{dzjowng} (Guangyun, in Baxter's transcription), glossed 'to go to; to follow'. 'From' is a later sense, assumed here, not quoted. It is the cóng of \emph{cónglái}.
\end{cousinweb}

\begin{grammarlens}[title={What you've built}]
\zh{从} comes before the starting point, and the whole phrase goes before the verb: \zh{我}, \zh{从中国}, \zh{来}.

\zh{到} after it marks where things end: \zh{从这儿到大学}.
\end{grammarlens}

\subsection*{Your turn}
\begin{itemize}
\raggedright
  \item \emph{Say it:} \emph{cóng}
  \item \emph{Say it:} \emph{cóng}, once more
  \item \emph{Say it:} \zh{从这儿到大学} — \emph{cóng zhèr dào dàxué} — from here to the university
  \item \emph{Say it:} \zh{我从中国来} — \emph{wǒ cóng Zhōngguó lái} — I come from China.
  \item \emph{Recall:} say \emph{cónglái}, then \zh{从}, and hear the same cóng
\end{itemize}

\begin{tcolorbox}[breakable,colback=teal!4,colframe=teal!35!black,title={Before you move on}]
What does \zh{从} mean? (\textquotedblleft{}From\textquotedblright{}.) Where does \zh{从} and its place go? (\textbf{Before the verb: \zh{我从中国来}}.)
\end{tcolorbox}
\section[\texorpdfstring{\zh{从} (cóng)}{从 (cóng)}]{\zh{从} — \zh{人} twice, the left one narrower}

\label{lesson:ZH-W155-cong-from}



\begin{quote}
The word first, then the character inside it.

\begin{itemize}
\raggedright
  \item \emph{Recall:} say \emph{cóng} — \textbf{R1}, one lesson back, before it decays
\end{itemize}

\begin{quote}
\zh{人} \zh{个}
\end{quote}

\emph{Read it:} each one aloud, with nothing said first — \emph{rén}, \emph{gè}
\end{quote}

\subsection*{Script — \zh{人} twice, the left one narrower}
\hlblockfigure{\detokenize{figures/ZH-W155-cong-from-filmstrip.pdf}}{How \zh{从} is written, stroke by stroke}

\begin{quote}
\zh{从}
\end{quote}

\textbf{One piece you own, written twice.}

\textbf{The left \zh{人}, two strokes}, narrow: the left-falling stroke, then a short stroke falling to the right that stops early, almost a dot, so the second \zh{人} has room.

\textbf{The right \zh{人}, two strokes}, full size: the left-falling stroke, then the long stroke falling to the right.

2 + 2 = four. One person following another.

\subsection*{Your turn}
\begin{itemize}
\raggedright
  \item \emph{Write it:} \zh{从}, all four, with the model in front of you
  \item \emph{Cover:} write it again from memory, and count the pieces, not the strokes
\end{itemize}

\begin{tcolorbox}[breakable,colback=teal!4,colframe=teal!35!black,title={Before you move on}]
How many strokes in \zh{从}? (\textbf{Four}.) Which \zh{人} is narrower? (\textbf{The left}.) What happens to its second stroke? (\textbf{It stops early, almost a dot}.)

Source: \href{https://raw.githubusercontent.com/chanind/hanzi-writer-data/68d10a4b21150cae5e1ebbd223eed289cf32d90c/data/\%E4\%BB\%8E.json}{Hanzi Writer Data \zh{从}, pinned stroke paths and medians}.
\end{tcolorbox}
\section[\texorpdfstring{\zh{从八点到十点} (cóng bā diǎn dào shí diǎn)}{从八点到十点 (cóng bā diǎn dào shí diǎn)}]{\zh{从八点到十点} — from and to}

\label{lesson:ZH-C155-cong-dao}



\begin{quote}
Two characters you write, a start and an end.

\begin{itemize}
\raggedright
  \item \emph{Read it:} \zh{从这儿}, aloud — \emph{cóng zhèr}
  \item \emph{Recall:} say \emph{jǐ diǎn}, what time
  \item \emph{Recall:} say \emph{zhǔnshí}, on time
\end{itemize}
\end{quote}

\begin{grammarlens}[title={Grammar Lens: \zh{从} … \zh{到} …, a start and an end}]
\begin{quote}
\textbf{\zh{从} marks where something starts, \zh{到} where it ends. A time or a place goes after each.}
\end{quote}

\begin{quote}
\textbf{\zh{从八点到十点}} \emph{cóng bā diǎn dào shí diǎn} — from eight to ten o'clock \textbf{\zh{从这儿到大学}} \emph{cóng zhèr dào dàxué} — from here to the university
\end{quote}

On a sign, \zh{到} often stands alone between two days or two times, with \zh{从} left out:

\begin{quote}
\textbf{\zh{星期一到星期五}} \emph{xīngqīyī dào xīngqīwǔ} — Monday to Friday
\end{quote}

In a sentence the whole phrase goes before the verb:

\begin{quote}
\textbf{\zh{我从八点到十点上课}} \emph{wǒ cóng bā diǎn dào shí diǎn shàngkè} — I have class from eight to ten.
\end{quote}
\end{grammarlens}

\subsection*{Your turn}
Say these aloud:

\begin{itemize}
\raggedright
  \item \zh{从八点到十点} — \emph{cóng bā diǎn dào shí diǎn} — from eight to ten o'clock
  \item \zh{从这儿到大学} — \emph{cóng zhèr dào dàxué} — from here to the university
  \item \zh{星期一到星期五} — \emph{xīngqīyī dào xīngqīwǔ} — Monday to Friday
  \item \zh{我从八点到十点上课} — \emph{wǒ cóng bā diǎn dào shí diǎn shàngkè} — I have class from eight to ten.
\end{itemize}

\begin{tcolorbox}[breakable,colback=teal!4,colframe=teal!35!black,title={Before you move on}]
What does \zh{从} mark? (\textbf{Where it starts}.) And \zh{到}? (\textbf{Where it ends}.) What can a sign leave out? (\textbf{\zh{从}: \zh{星期一到星期五}}.)
\end{tcolorbox}
\section[chūfā]{chūfā — to set off; to depart}

\label{lesson:ZH-C155-chufa}



\begin{quote}
Before the new one: read the last lesson's phrase, then going out of the door.

\begin{itemize}
\raggedright
  \item \emph{Read it:} \zh{从八点到十点}, aloud — \emph{cóng bā diǎn dào shí diǎn}
  \item \emph{Recall:} say \emph{chūmén}, to go out — \textbf{R3}, durable
\end{itemize}
\end{quote}

\subsection*{You'll want to know: chūfā}
\emph{chūfā} — \textquotedblleft{}to set off; to depart\textquotedblright{}.

\begin{quote}
\emph{huǒchē bā diǎn chūfā} — The train leaves at eight.
\end{quote}

\begin{quote}
\emph{wǒmen jǐ diǎn chūfā} — What time do we set off?
\end{quote}

Its second character comes later; for now it is a word for the ear.

\begin{cousinweb}
\emph{chūfā} is built from \zh{出} 'to go out', the \zh{出} you write, and fā 'to send off; to start', the fā of \emph{fāshāo}. Both are inherited from Middle Chinese, \emph{tsyhwit} and \emph{pjot} (Guangyun, in Baxter's transcription), glossed 'to go forward; to appear' and 'to start up'.
\end{cousinweb}

\begin{grammarlens}[title={What you've built}]
\emph{chūfā} says when a journey starts; \zh{到} says when it arrives: \emph{bā diǎn chūfā, shí diǎn dào}.

On a timetable the start is often written \zh{开}, as a train or a bus \zh{开} when it leaves.
\end{grammarlens}

\subsection*{Your turn}
\begin{itemize}
\raggedright
  \item \emph{Say it:} \emph{chūfā}
  \item \emph{Say it:} \emph{chūfā}, once more
  \item \emph{Say it:} \emph{huǒchē bā diǎn chūfā} — The train leaves at eight.
  \item \emph{Say it:} \emph{wǒmen jǐ diǎn chūfā} — What time do we set off?
  \item \emph{Recall:} say \emph{chūmén}, then \emph{chūfā}
\end{itemize}

\begin{tcolorbox}[breakable,colback=teal!4,colframe=teal!35!black,title={Before you move on}]
What does \emph{chūfā} mean? (\textquotedblleft{}To set off\textquotedblright{}.) And to arrive? (\textbf{\zh{到}}.)
\end{tcolorbox}
\section[\texorpdfstring{\zh{开门} (kāimén)}{开门 (kāimén)}]{\zh{开门} (kāimén) — to open; to be open}

\label{lesson:ZH-C155-kaimen}



\begin{quote}
Before the new one: last lesson's word, then the two words this one is made of.

\begin{itemize}
\raggedright
  \item \emph{Recall:} say \emph{chūfā}, to set off — \textbf{R1}, one lesson back
  \item \emph{Recall:} say \emph{kāi}, to open — \textbf{R3}, durable
  \item \emph{Recall:} say \zh{门} — \emph{mén}, a door — \textbf{R3}
\end{itemize}
\end{quote}

\subsection*{You'll want to know: \zh{开门}}
\textbf{\zh{开门}} — \emph{kāimén} — \textquotedblleft{}to open (for business); to open the door\textquotedblright{}.

\begin{quote}
\textbf{\zh{星期天不开门}} \emph{xīngqītiān bù kāimén} — It's closed on Sundays.
\end{quote}

\begin{quote}
\textbf{\zh{早上九点开门}} \emph{zǎoshang jiǔ diǎn kāimén} — It opens at nine in the morning.
\end{quote}

\begin{cousinweb}
\textbf{\zh{开门}} (\emph{kāimén}) is built from \zh{开} 'to open' and \zh{门} 'a door'. Both are inherited from Middle Chinese, \emph{khoj} and \emph{mwon} (Guangyun, in Baxter's transcription), \emph{khoj} glossed 'to open; to undo'.
\end{cousinweb}

\begin{grammarlens}[title={What you've built}]
A shop that \zh{开门} is open for business, not only its door. On a sign the hours come first: \zh{早上九点到晚上八点开门}.

\zh{不} stays \emph{bù} before the level tone of kāi: \zh{不开门}.
\end{grammarlens}

\subsection*{Your turn}
\begin{itemize}
\raggedright
  \item \emph{Say it:} \emph{kāimén}
  \item \emph{Say it:} \emph{kāimén}, once more
  \item \emph{Say it:} \zh{星期天不开门} — \emph{xīngqītiān bù kāimén} — It's closed on Sundays.
  \item \emph{Say it:} \zh{早上九点开门} — \emph{zǎoshang jiǔ diǎn kāimén} — It opens at nine in the morning.
  \item \emph{Recall:} say \emph{kāi} and \zh{门}, then \zh{开门}
\end{itemize}

\begin{tcolorbox}[breakable,colback=teal!4,colframe=teal!35!black,title={Before you move on}]
What does \zh{开门} mean on a shop? (\textbf{It is open for business}.) Say it is closed on Sundays. (\textbf{\zh{星期天不开门}}.)
\end{tcolorbox}
\section[guānmén]{guānmén — to close; to shut the door}

\label{lesson:ZH-C155-guanmen}



\begin{quote}
Before the new one: last lesson's word, then the verb for shutting.

\begin{itemize}
\raggedright
  \item \emph{Recall:} say \zh{开门} — \emph{kāimén} — \textbf{R1}, one lesson back
  \item \emph{Recall:} say \emph{guān}, to close — \textbf{R3}, durable
\end{itemize}
\end{quote}

\subsection*{You'll want to know: guānmén}
\emph{guānmén} — \textquotedblleft{}to close (for the day); to shut the door\textquotedblright{}.

\begin{quote}
\emph{wǎnshang bā diǎn guānmén} — It closes at eight in the evening.
\end{quote}

Its first character comes later; for now it is a word for the ear.

\begin{cousinweb}
\emph{guānmén} is built from \emph{guān} 'to close', the word you know, and \zh{门} 'a door'. \emph{guān} is inherited from Middle Chinese \emph{kwaen} (Guangyun, in Baxter's transcription), which the dictionary explains, through an older book, as a wooden bar held across a door.
\end{cousinweb}

\begin{grammarlens}[title={What you've built}]
\zh{开门} and \emph{guānmén} bracket a shop's day: \emph{zǎoshang jiǔ diǎn kāimén, wǎnshang bā diǎn guānmén}.

A sign often gives only the span, with \zh{到} between: \zh{九点到八点}.
\end{grammarlens}

\subsection*{Your turn}
\begin{itemize}
\raggedright
  \item \emph{Say it:} \emph{guānmén}
  \item \emph{Say it:} \emph{guānmén}, once more
  \item \emph{Say it:} \emph{wǎnshang bā diǎn guānmén} — It closes at eight in the evening.
  \item \emph{Recall:} say \zh{开门}, then \emph{guānmén}
\end{itemize}

\begin{tcolorbox}[breakable,colback=teal!4,colframe=teal!35!black,title={Before you move on}]
What does \emph{guānmén} mean? (\textquotedblleft{}To close for the day\textquotedblright{}.) And to open? (\textbf{\zh{开门}}.)
\end{tcolorbox}
\section[(reading)]{(reading) — a timetable and opening hours}

\label{lesson:ZH-C155-reading}



\begin{quote}
Two notices, the kind you read standing up: a station timetable, then a bookshop's door.

\begin{itemize}
\raggedright
  \item \emph{Recall:} say \emph{guānmén}, to close — \textbf{R1}, one lesson back
  \item \emph{Recall:} say \emph{fēnzhōng}, a minute — \textbf{R2}, the chapter on the clock
\end{itemize}

\begin{quote}
\zh{十点半}
\end{quote}

\emph{Read it:} it aloud, with nothing said first — \emph{shí diǎn bàn}
\end{quote}

\begin{tcolorbox}[breakable,skin=enhanced,colback=orange!6,colframe=orange!45!black,boxrule=0.5pt,arc=1mm,left=6pt,right=6pt,top=4pt,bottom=4pt,fonttitle=\bfseries,title={Read this}]
\begin{quote}
\zh{从这儿到大学的火车} \zh{早上八点开} \zh{十点半到} \zh{下午两点开} \zh{四点半到} \zh{晚上七点开} \zh{九点到}
\end{quote}

\begin{quote}
\zh{书} \zh{星期一到星期五} \zh{早上九点到晚上八点开门} \zh{星期六} \zh{上午十点到下午五点开门} \zh{星期天不开门}
\end{quote}

\emph{Read it:} both notices through once without stopping; do not translate as you go

\emph{Read it:} the timetable again, and find the train that arrives before eleven
\end{tcolorbox}

\begin{tcolorbox}[breakable,skin=enhanced,colback=orange!6,colframe=orange!45!black,boxrule=0.5pt,arc=1mm,left=6pt,right=6pt,top=4pt,bottom=4pt,fonttitle=\bfseries,title={Read this}]
\begin{quote}
\emph{cóng zhèr dào dàxué de huǒchē} \emph{zǎoshang bā diǎn kāi, shí diǎn bàn dào} \emph{xiàwǔ liǎng diǎn kāi, sì diǎn bàn dào} \emph{wǎnshang qī diǎn kāi, jiǔ diǎn dào}
\end{quote}

\begin{quote}
\emph{shū} \emph{xīngqīyī dào xīngqīwǔ} \emph{zǎoshang jiǔ diǎn dào wǎnshang bā diǎn kāimén} \emph{xīngqīliù} \emph{shàngwǔ shí diǎn dào xiàwǔ wǔ diǎn kāimén} \emph{xīngqītiān bù kāimén}
\end{quote}

A timetable leaves out every word it can. \zh{开} is the train leaving and \zh{到} is it arriving; the shop's name is one character, \zh{书}, and its hours are two times with \zh{到} between them and \zh{开门} after.
\end{tcolorbox}

\subsection*{Your turn}
\begin{itemize}
\raggedright
  \item You must be at the university by eleven: which train? (\textbf{The eight o'clock: \zh{早上八点开} \zh{十点半到}}.)
  \item Can you buy a book on Sunday? (\textbf{No: \zh{星期天不开门}}.)
  \item What time does the bookshop open on Saturday? (\textbf{Ten in the morning: \zh{上午十点}}.)
\end{itemize}

\begin{itemize}
\raggedright
  \item \emph{Say it:} the timetable aloud, once, without stopping
\end{itemize}

\begin{tcolorbox}[breakable,colback=teal!4,colframe=teal!35!black,title={Before you move on}]
When does the last train arrive? (\textbf{\zh{九点}}, nine in the evening.) Which day is the shop shut? (\textbf{\zh{星期天}}.) How late is it open on a weekday? (\textbf{\zh{晚上八点}}.)
\end{tcolorbox}
\section[Practice]{\zh{火车十点半到}}

\label{lesson:ZH-C155-practice}



\begin{quote}
The chapter closes here: the same timetable and the same door, asked about aloud.

\begin{itemize}
\raggedright
  \item \emph{Recall:} say \zh{星期天不开门} — \emph{xīngqītiān bù kāimén} — \textbf{R1}, the reading
  \item \emph{Recall:} say \emph{chūfā}, to set off — \textbf{R2}
  \item \emph{Recall:} say \emph{guānmén}, to close — \textbf{R2}
\end{itemize}
\end{quote}

\subsection*{The exchange}
\begin{quote}
\textbf{A:} \emph{qǐngwèn, huǒchē jǐ diǎn chūfā} — Excuse me, what time does the train leave?
\end{quote}

\begin{quote}
\textbf{B:} \zh{早上八点开}
\end{quote}

\begin{quote}
\emph{zǎoshang bā diǎn kāi} — At eight in the morning.
\end{quote}

\begin{quote}
\textbf{A:} \emph{jǐ diǎn dào dàxué} — What time does it get to the university?
\end{quote}

\begin{quote}
\textbf{B:} \zh{十点半到}
\end{quote}

\begin{quote}
\emph{shí diǎn bàn dào} — At half past ten.
\end{quote}

\begin{quote}
\textbf{A:} \zh{书} \zh{星期天开门吗}
\end{quote}

\begin{quote}
\emph{shū, xīngqītiān kāimén ma} — The bookshop: is it open on Sunday?
\end{quote}

\begin{quote}
\textbf{B:} \zh{星期天不开门}
\end{quote}

\begin{quote}
\emph{xīngqītiān bù kāimén} — It's closed on Sundays.
\end{quote}

\begin{quote}
\textbf{A:} \zh{星期六呢}
\end{quote}

\begin{quote}
\emph{xīngqīliù ne} — And Saturday?
\end{quote}

\begin{quote}
\textbf{B:} \zh{上午十点到下午五点开门}
\end{quote}

\begin{quote}
\emph{shàngwǔ shí diǎn dào xiàwǔ wǔ diǎn kāimén} — Open from ten to five.
\end{quote}

\begin{quote}
\textbf{A:} \emph{xiàwǔ wǔ diǎn guānmén} — So it closes at five.
\end{quote}

The questions come in pinyin, because \emph{jǐ} is not a character you write yet; the answers come off the notices, in characters.

\begin{tcolorbox}[breakable,skin=enhanced,colback=orange!6,colframe=orange!45!black,boxrule=0.5pt,arc=1mm,left=6pt,right=6pt,top=4pt,bottom=4pt,fonttitle=\bfseries,title={Read this}]
\emph{Read it:} the exchange through once without stopping, then answer the question below
\end{tcolorbox}

\begin{grammarlens}[title={What you've built — reading the times off a notice}]
A start, an end, and a verb for what happens:

\begin{itemize}
\raggedright
  \item \textbf{\zh{早上八点开} \zh{十点半到}} \emph{zǎoshang bā diǎn kāi, shí diǎn bàn dào} — leaves at eight, arrives at half past ten.
  \item \textbf{\zh{从八点到十点}} \emph{cóng bā diǎn dào shí diǎn} — from eight to ten.
  \item \textbf{\zh{星期一到星期五}} \emph{xīngqīyī dào xīngqīwǔ} — Monday to Friday.
  \item \emph{wǎnshang bā diǎn guānmén} — closes at eight in the evening.
\end{itemize}

\textbf{Listen.} \emph{Huǒchē jǐ diǎn dào?}

\textbf{Speak.} Give the hours of a shop you know, from opening to closing.

\emph{Write it:} \zh{十点半到} — then \zh{从}, one \zh{人} following another
\end{grammarlens}

\subsection*{Your turn}
\begin{itemize}
\raggedright
  \item \emph{Say it:} \emph{huǒchē jǐ diǎn dào} — What time does the train arrive?
  \item \emph{Say it:} \zh{十点半到} — \emph{shí diǎn bàn dào} — At half past ten.
  \item \emph{Say it:} the whole exchange, both parts
  \item \emph{Write it:} \zh{从八点到十点}
\end{itemize}

\begin{tcolorbox}[breakable,colback=teal!4,colframe=teal!35!black,title={Before you move on}]
Say the train arrives at half past ten. (\textbf{\zh{十点半到}}.) Say it's closed on Sundays. (\textbf{\zh{星期天不开门}}.) Say from eight to ten. (\textbf{\zh{从八点到十点}}.)
\end{tcolorbox}
